% Algorithm Documentation Requirements (MAIN POINT 1). Built by the future pass
% from the instructions below. Do not process now. The section both states what
% algorithm documentation the bill requires to be filed and demonstrates that
% requirement against the recorded VVUQ-02 codegen and execution artifacts. The
% 1790-line comparison.json and the 1001-row sample_h2_sensor.csv must be
% featured here.

[Write this as connected legislative prose plus two tables. Movement 1 states the
documentation duty the bill imposes; Movement 2 demonstrates that the duty is
already met in practice by the recorded VVUQ-02 artifacts, which is the proof the
requirement is feasible rather than burdensome. Keep file names in the tables.]

[Movement 1, The documentation duty. Specify what a sponsor must document and file
for each algorithm that governs a robot-patient interaction, framed so the
verification record precedes the generated code (rigor above the code itself).
The required filing items, to be presented as a numbered list and a table, are: a
plain-language summary of the algorithm's function; the decision-logic or
architecture description; a description of the data sets used to train or
condition the algorithm; the VVUQ gate bindings (which external standard governs
each gate and the numeric threshold enforced); the verification-before-generation
record (the gate result that cleared before the code was generated or executed);
and the determinism and reproducibility statement (seed and byte-for-byte
regeneration). Align these items to the state utilization-review precedents so the
bill is consistent with them: the Texas AI compliance statement of summary,
decision-logic tree, training-data description, and attestation
(\texttt{tx-sb1822}); the New York algorithm-and-training-data submission
(\texttt{ny-a9149}); and the California individualized-data and transparency duty
(\texttt{ca-sb1120}). Anchor the device-software framing to \texttt{cfr-devices}
and the credibility framing to \texttt{fda2023credibility} and \texttt{asmevv40}.]

[Movement 2, Demonstration from the recorded artifacts. Show that the VVUQ-02
codegen and execution records already produce exactly this documentation, drawing
from \texttt{VVUQ-02/codegen/} and \texttt{VVUQ-02/execution/}. Describe the
generated platform and its documentation set from the docs in
\texttt{VVUQ-02/codegen/docs/} (architecture, h2 robot specification, h2 hand
kinematics, sensor spec, perception stack, autonomy policy, the gate spec, and the
standards map) and the frozen configuration in \texttt{VVUQ-02/codegen/config/}.
Two artifacts must be featured by name and described in detail, each in its own
table, with the parent path given once at the top.]

[Featured artifact 1, the four-entrant tournament. File:
\texttt{papers/VVUQ-02/codegen/results/comparison.json}, 1790 lines, the result
of a substantial code generation processing feat and its subsequent execution.
Describe it as the machine-readable record of a 32-iteration sweep at seed
20260525 with four rounds per iteration, comparing four entrants
(\texttt{H2\_Surgical\_1\_0}, the single mobile humanoid; \texttt{PancreSpeed\_1\_0},
the eight-arm cart baseline; \texttt{da\_Vinci\_successor\_2030}, a teleoperated
telemanipulator; and \texttt{Dutch\_human\_baseline}, a 2025 human-surgeon
cohort), each round carrying a composite score for each side, a winner, a
confidence value, and explicit caveats that always include the simulation-against-
simulation qualifier. Report the headline numbers from the record (composite range
near 67.584 to 94.603; the humanoid mean near 93.334 placing second to PancreSpeed
near 93.782 by under half a composite point) and state that the file regenerates
byte for byte from the seed. Present this in one table; do not transcribe the
JSON.]

[Featured artifact 2, the humanoid sensor stream. File:
\texttt{papers/VVUQ-02/codegen/data/sample\_h2\_sensor.csv}, 1001 lines (1 header
plus 1000 non-repetitive data rows) across 27 columns, the flattened publication
sample of the robot positional data produced by the code generation and confirmed
in execution. Describe the columns by group (tick, hand identifier, phase, grasp
state; the seven per-arm joint angles \texttt{q\_arm\_1} through \texttt{q\_arm\_7};
the five per-finger forces \texttt{finger\_force\_1} through
\texttt{finger\_force\_5}; the end-effector position \texttt{ee\_pos\_x/y/z} in the
patient frame; and the vessel-proximity, ring-tension, balance, e-stop, and
collision channels), so the reader sees that the data covers both hands, all five
fingers, the arms, and the balance state with no repeated rows, deterministic from
seed 20260525. Note the companion \texttt{sample\_h2\_sensor.jsonl} and the
human-surgeon baseline \texttt{human\_surgeon\_baseline.csv}. Present the column
groups in one table; do not transcribe the CSV.]

[Close Movement 2 by connecting the demonstration back to the duty: because these
artifacts were generated and executed autonomously and reproducibly, the bill's
documentation requirement is shown to be practical and inexpensive. Defer the
honest account of what could not be run or had to be approximated to
\texttt{sections/limitations\_future.tex}, and the not-physically-attached
treatment of these same files to \texttt{sections/supporting\_documentation.tex}.
Cite \texttt{kawchak2026vvuq02} and \texttt{repo-cancer-automated}.]
